\section*{Capítulo \ref{ch:b2:biomolecules} --- Biomoléculas: aminoácidos, açúcares, lipídios e nucleotídeos}

\begin{solution}{exo:b2:biomolecules:1}
\begin{itemize}
  \item Entre pH 3 e 9, o \omterm{def:b2:biomolecules:amino-acid}{zwitteríon} \ce{H3N+-CH(CH3)-COO-}.
  \item A pH 1, o cátion \ce{H3N+-CH(CH3)-COOH}.
  \item A pH 12, o ânion \ce{H2N-CH(CH3)-COO-}.
\end{itemize}
\end{solution}

\begin{solution}{exo:b2:biomolecules:2}
Glicina: $(2.35 + 9.78)/2 \approx 6.1$. Alanina: $(2.35 + 9.87)/2 \approx 6.1$.
% ledger: pka:glycine1, pka:glycine2, pka:alanine1, pka:alanine2
\end{solution}

\begin{solution}{exo:b2:biomolecules:3}
\ce{H2N-CH(CH3)-CO-NH-CH2-COOH}: a alanina fornece o N-terminal, a glicina o C-terminal. Gly--Ala,
\ce{H2N-CH2-CO-NH-CH(CH3)-COOH}, é um isômero constitucional: o grupo amino livre está na glicina.
\end{solution}

\begin{solution}{exo:b2:biomolecules:4}
\ce{NH2} à esquerda: L. Prioridades \ce{NH2} > \ce{COOH} > \ce{CH2OH} > \ce{H}: a L-serina é
(\textit{S}).
\end{solution}

\begin{solution}{exo:b2:biomolecules:5}
Grupos: amônio N-terminal (cerca de 9), \omterm{def:b2:biomolecules:amino-acid}{cadeia lateral} da lisina (10.7), \omterm{def:b2:biomolecules:amino-acid}{cadeia lateral} do aspartato
(3.9), grupo carboxílico C-terminal (cerca de 2). pH 1: $+1 +1 + 0 + 0 = +2$. pH 7: $+1 +1 -1 -1 = 0$. pH 12:
$0 + 0 - 1 - 1 = -2$ (a \omterm{def:b2:biomolecules:amino-acid}{cadeia lateral} da lisina, 1.3 unidade abaixo de pH 12, está sobretudo neutra).
\end{solution}

\begin{solution}{exo:b2:biomolecules:6}
O cloreto de tionila e o aquecimento atacariam também outros grupos e, sobretudo, o \omterm{def:b2:acyl-substitution:derivative}{cloreto de acila} muito
reativo de um aminoácido perde facilmente sua configuração (o hidrogênio $\alpha$ torna-se ácido): o
\omterm{def:b2:biomolecules:peptide}{peptídeo} seria parcialmente racemizado. Um \omterm{def:b2:biomolecules:peptide}{agente de acoplamento} ativa o ácido à temperatura ambiente,
em condições brandas.
\end{solution}

\begin{solution}{exo:b2:biomolecules:7}
O5 ligado a C2: um anel de C2, C3, C4, C5 e O, cinco átomos (uma furanose). Em \omterm{def:b2:biomolecules:anomer}{projeção de Haworth}, O do
anel atrás, C2 à direita com \ce{CH2OH} (C1) para baixo e \ce{OH} para cima ($\beta$); \ce{OH} de C3 para
cima (à esquerda em Fischer), \ce{OH} de C4 para baixo, C5 levando \ce{CH2OH} (C6) para cima.
\end{solution}

\begin{solution}{exo:b2:biomolecules:8}
$x_\alpha = (80.2 - 52.8)/(150.7 - 52.8) = 27.4/97.9 \approx 0.28$: 28~\% de $\alpha$, 72~\% de $\beta$.
\end{solution}

\begin{solution}{exo:b2:biomolecules:9}
A maltose conserva um hemiacetal livre em sua segunda glicose, que se abre em aldeído: redutora. Na
sacarose, os dois \omterm{def:b2:biomolecules:anomer}{carbonos anoméricos} estão engajados na \omterm{def:b2:biomolecules:glycoside}{ligação glicosídica}: nenhum hemiacetal, não
redutora. O ácido diluído hidrolisa a \omterm{def:b2:biomolecules:glycoside}{ligação glicosídica}, um acetal: a glicose e a frutose são
liberadas, e ambas reduzem.
\end{solution}

\begin{solution}{exo:b2:biomolecules:10}
Trioleína, \qty{885.45}{g/mol}, consome três \ce{KOH} (\qty{56.105}{g/mol}): $3 \times 56.105/885.45 =
0.190$~g por g, um índice de \omterm{def:b2:acyl-substitution:saponification}{saponificação} de 190. Cadeias mais curtas significam uma massa molar menor
para os mesmos três grupos éster: mais mols de gordura por grama, mais KOH.
% ledger: aw:C, aw:H, aw:O, aw:K
\end{solution}

\begin{solution}{exo:b2:biomolecules:11}
Proteja a amina da glicina como carbamato de \textit{terc}-butoxicarbonila, e o ácido da alanina como
éster metílico. Acople com uma carbodi-imida. Remova o carbamato com ácido e o éster com base diluída
(condições ortogonais, de modo que cada extremidade pode ser liberada sozinha): Gly--Ala.
\end{solution}

\begin{solution}{exo:b2:biomolecules:12}
5$'$-ACGCAT-3$'$ (lida de modo antiparalelo: A se pareia com T, G com C). Pares: A--T, T--A, G--C, C--G, G--C, T--A:
$3 \times 2 + 3 \times 3 = 15$ ligações de hidrogênio.
\end{solution}

\begin{solution}{pb:b2:biomolecules:1}
\textbf{1.} Ácido L-aspártico, L-fenilalanina e metanol.
\textbf{2.} Nos dois, \ce{COOH} no alto, \omterm{def:b2:biomolecules:amino-acid}{cadeia lateral} embaixo (\ce{CH2COOH}, \ce{CH2C6H5}), \ce{NH2}
à esquerda.
\textbf{3.} (\textit{S}) para ambos.
\textbf{4.} Dois estereocentros: quatro estereoisômeros.
\textbf{5.} \ce{H2N-CH(CH2COOH)-CO-NH-CH(CH2C6H5)-COOCH3}: a \omterm{def:b2:biomolecules:peptide}{ligação peptídica} une os dois
resíduos; o éster é o \ce{COOCH3} do C-terminal.
\textbf{6.} Os receptores do paladar são quirais: eles ligam um estereoisômero e não sua imagem
especular nem seus diastereoisômeros, como dizia a introdução.
\textbf{7.} O \ce{NH3+} N-terminal e o \ce{COOH} da \omterm{def:b2:biomolecules:amino-acid}{cadeia lateral} do ácido aspártico; o ácido C-terminal
é um éster metílico, sem hidrogênio ácido.
\textbf{8.} A pH 2: \ce{NH3+} ($+1$), \ce{COOH} da \omterm{def:b2:biomolecules:amino-acid}{cadeia lateral} (0): $+1$.
\textbf{9.} A pH 7: $+1 - 1 = 0$.
\textbf{10.} A pH 12: $0 - 1 = -1$.
\textbf{11.} Entre os dois $\mathrm pK_a$ dos grupos que mudam: $(3.9 + 9.9)/2 \approx 6.9$.
\textbf{12.} No dipeptídeo, o carbono da amina leva uma amida em vez de um carboxilato: os grupos e as
cargas vizinhas mudam, e os $\mathrm pK_a$ também (o amônio N-terminal de um \omterm{def:b2:biomolecules:peptide}{peptídeo} é nitidamente
mais ácido que o de um aminoácido livre). O modelo dá a forma, não os valores exatos.
\textbf{13.} O ácido e a amina formam primeiro um sal; aquecer daria misturas (Asp--Asp, cadeias mais
longas, o grupo carboxílico errado) e racemizaria os estereocentros.
\textbf{14.} Seu grupo amino e seu grupo carboxílico da \omterm{def:b2:biomolecules:amino-acid}{cadeia lateral}.
\textbf{15.} Ele transforma o \ce{OH} do ácido em um bom grupo de saída, um derivado ativado que a amina
ataca à temperatura ambiente (substituição acílica).
\textbf{16.} O $\beta$-dipeptídeo, ligado pela \omterm{def:b2:biomolecules:amino-acid}{cadeia lateral}: um isômero do aspartame, não o aspartame.
\textbf{17.} A amina como carbamato de benziloxicarbonila e a \omterm{def:b2:biomolecules:amino-acid}{cadeia lateral} como éster benzílico: os
dois são removidos juntos por hidrogenólise sobre paládio, que deixa intacto o éster metílico
(ortogonal a ele).
\textbf{18.} Um ácido ativado tem um hidrogênio $\alpha$ ácido; a base o retira e o estereocentro
racemiza.
\textbf{19.} Proteja o ácido aspártico (carbamato no N, éster benzílico na \omterm{def:b2:biomolecules:amino-acid}{cadeia lateral}); ative o
$\alpha$-\ce{COOH} livre com o \omterm{def:b2:biomolecules:peptide}{agente de acoplamento}; acrescente o éster metílico da L-fenilalanina;
faça a hidrogenólise.
\textbf{20.} O éster metílico e a \omterm{def:b2:biomolecules:peptide}{ligação peptídica}; o éster primeiro, pois uma amida é muito menos reativa.
\textbf{21.} O dipeptídeo Asp--Phe (ácido livre) e o metanol.
\textbf{22.} A molécula doce é consumida: seus produtos de hidrólise não a substituem.
\textbf{23.} $14 \times 12.011 + 18 \times 1.008 + 2 \times 14.007 + 5 \times 15.999 =
\qty{294.307}{g/mol}$.
\textbf{24.} $0.180/294.307 = \qty{6.116e-4}{mol}$, \qty{0.612}{mmol}.
\textbf{25.} Um metanol por aspartame: $0.6116 \times 32.042 = \boldsymbol{\approx \qty{19.6}{mg}}$ de
metanol.
\end{solution}
